% Generated draft: requires mathematical and source review.
\documentclass[11pt]{article}
\usepackage[margin=25mm]{geometry}
\usepackage{fontspec}
\setmainfont{Latin Modern Roman}
\usepackage{amsmath,amssymb,mathtools,unicode-math}
\setmathfont{Latin Modern Math}
\usepackage{enumitem,needspace,xurl,hyperref}
\setlength{\parindent}{0pt}
\setlength{\parskip}{5pt}
\setlength{\emergencystretch}{3em}
\newcommand{\N}{\mathbb{N}}
\newcommand{\Z}{\mathbb{Z}}
\newcommand{\Q}{\mathbb{Q}}
\newcommand{\R}{\mathbb{R}}
\newcommand{\C}{\mathbb{C}}
\newcommand{\F}{\mathbb{F}}
\newcommand{\A}{\mathbb{A}}
\newcommand{\PP}{\mathbb{P}}
\newcommand{\E}{\mathbb{E}}
\newcommand{\eps}{\varepsilon}
\newcommand{\dd}{\,\mathrm{d}}
\begin{document}
\section*{UnsolvedMath ID: 30002283}
\textbf{Source problem code:} catalogue entry
\subsection{Definitions and mathematical statement}
A del Pezzo surface over a field \(k\) is a smooth projective geometrically integral surface \(S\) whose anticanonical divisor \(-K_S\) is ample; its degree is \(K_S^2\). The supplied background concerns quartic del Pezzo fibrations, whose smooth geometric fibres have degree \(4\), whereas the extracted question says only ``del Pezzo fibrations.'' The intended degree, ground field, and allowed singular fibres therefore require confirmation.

Write \(f\colon X\to B\) for a fibration of the intended class. The record does not define its height, its discriminant, or the monodromy representation. Denote the intended height provisionally by \(h(f)\). If the discriminant is an effective divisor on a smooth base curve, ``square-free'' means that every component has multiplicity one; the record does not specify the discriminant construction or the base for the general questions. ``Maximal monodromy'' requires a specified representation and ambient group, neither of which is supplied.

Subject to these unresolved conventions, the source asks the following questions.
\begin{enumerate}
\item Does there exist a del Pezzo fibration with square-free discriminant and \(h(f)=2\)? Does there exist one with \(h(f)=6\) and maximal monodromy? The wording does not determine whether the square-free-discriminant condition also applies to the height-six question.
\item For a fixed height \(h\), is the parameter or moduli space of the intended fibrations of height \(h\) irreducible? The record does not identify this space or its equivalence relation. This formulation interprets ``are those of a given height irreducible'' as a question about a moduli space; that interpretation is unconfirmed.
\item How can the Prym varieties associated with these fibrations be characterized? Neither the associated covering construction nor the desired kind of characterization is specified. For a finite morphism \(\pi\colon C'\to C\) of smooth projective curves, the usual Prym construction is
\[
\operatorname{Prym}(C'/C)=\bigl(\ker(\operatorname{Nm}_{\pi}\colon\operatorname{Jac}(C')\to\operatorname{Jac}(C))\bigr)^0,
\]
where the superscript denotes the identity component. This recalls a convention only; the record does not identify \(C'\), \(C\), \(\pi\), or an intended polarization.
\item For the particular height-twelve example \(f\colon X\to\mathbb P^1\) mentioned in the source, with its designated canonical section \(s_{\mathrm{can}}\), does there exist a section \(s\colon\mathbb P^1\to X\) such that
\[
f\circ s=\operatorname{id}_{\mathbb P^1},\qquad s\ne s_{\mathrm{can}}?
\]
The record supplies neither the construction of this particular \(X\) nor the definition of its canonical section, and does not specify the field over which sections are sought.
\end{enumerate}
These are source questions with unresolved formulation data, not assertions of existence, irreducibility, or a Prym characterization.
\subsection{Short English statement}
The source asks about existence at heights two and six, irreducibility at fixed height, associated Prym varieties, and additional sections of a particular height-twelve fibration over \(\mathbb P^1\). Definitions of height and monodromy, the relevant moduli space and Prym construction, and the particular height-twelve example are missing. The background suggests quartic del Pezzo fibrations, but the extracted question does not explicitly impose degree four.
\subsection{Sources}
\begin{itemize}
\item Supplied local UnsolvedMath registry record: permanent ID \(30002283\), exact source problem code \texttt{catalogue entry}, title ``Classification Problems for Del Pezzo Fibrations,'' in the Oberwolfach Reports Open Problems collection. The supplied proposed year is 2013; no individual proposer is identified.
\item Supplied source citation: \emph{Moduli Spaces in Algebraic Geometry} (2013), Oberwolfach Reports, pp.\ 343--392, DOI \texttt{10.4171/owr/2013/06}. \url{https://doi.org/10.4171/owr/2013/06}. The record attributes the questions to workshop open-problem discussions, without a precise internal page or speaker attribution.
\item Supplied literature source: Hassett--Tschinkel, \emph{Quartic del Pezzo surfaces}. \url{https://arxiv.org/abs/1301.7270}.
\item Supplied literature source: Hassett--Kresch--Tschinkel, \emph{Moduli of quartic del Pezzo fibrations}. \url{https://math.nyu.edu/~tschinke/papers/yuri/13noheighttwo/noheighttwo16.pdf}.
\item Literature assessment dated exactly \texttt{2026-08-21}: the supplied curation labels the problem open and states that later work classifies many heights and components of quartic del Pezzo fibrations, while the extracted multipart question, especially the Prym-variety characterization, is not fully resolved. This is supplied dated curation, not an independently verified literature assessment; the background explicitly requests independent verification of time-sensitive status and cited claims.
\item The supplied statement status is \texttt{corrected\_verified}. Its accompanying verification note reports that the mathematical content was checked against the official Oberwolfach report and its notation, context, and question form were normalized without changing the source-level claim. This records the supplied verification claim only; no source was independently consulted for this draft.
\end{itemize}
\end{document}
